% Audit de mod_stage : revue statique du code (contrôle d'accès, échappement et SQL, circuit du
% stage et intégrité des données, imports et exports, RGPD et sauvegarde, tests et maintenabilité).
% Compilation : latexmk -pdf audit-du-module.tex
\documentclass[11pt,a4paper]{article}
\usepackage[french]{babel}
\usepackage[T1]{fontenc}
\usepackage{lmodern,microtype,booktabs,tabularx,xcolor,hyperref,enumitem}
\usepackage[margin=2.2cm]{geometry}
\usepackage{url}
\expandafter\def\expandafter\UrlBreaks\expandafter{\UrlBreaks\do\a\do\b\do\c\do\d\do\e\do\f\do\g\do\h\do\i\do\j\do\k\do\l\do\m\do\n\do\o\do\p\do\q\do\r\do\s\do\t\do\u\do\v\do\w\do\x\do\y\do\z}
\DeclareUrlCommand\code{\urlstyle{tt}}
\Urlmuskip=0mu plus 2mu
\newcommand{\grav}[1]{\textbf{\textcolor{#1}{\textbullet}}}
\newcommand{\maj}{\grav{orange}}
\newcommand{\mineure}{\grav{gray}}
\newcommand{\crit}{\grav{red}}
\setlist[itemize]{itemsep=2pt,topsep=3pt}
\title{Audit de \texttt{mod\_stage} : revue du code}
\author{Sébastien Lefebvre}
\date{7 novembre 2026}

\begin{document}
\maketitle

\section{Méthode}

Les deux audits précédents ont mis le module à l'épreuve en situation : test de charge
(\code{test-de-charge.tex}) puis test de concurrence dans un Moodle réel
(\code{test-de-concurrence.tex}). Celui-ci est une revue statique du code, telle qu'on la ferait
avant une mise en production ou une publication dans le répertoire des plugins Moodle. Il porte
sur la branche \code{main} au commit \code{631e02b} (31~octobre 2026), version
\code{2026103101} (0.5.1) de \code{mod_stage} et \code{2026100200} (0.2.1) de
\code{mod_stagesynthesis}.

\begin{itemize}
  \item \textbf{Périmètre.} Les 55 scripts de page de \code{mod/stage}, \code{locallib.php}
    (8\,565 lignes, 259 fonctions), \code{lib.php}, les classes (formulaires, importateurs,
    PDF, événements, tâches, services web, confidentialité), le schéma (\code{install.xml},
    \code{upgrade.php}), la sauvegarde et la restauration, les gabarits Mustache, les fichiers de
    langue, les tests, la CI et la documentation ; soit 31\,000 lignes de PHP hors tests et
    bibliothèque tierce. \code{mod_stagesynthesis} n'est regardé que pour ce qui lui est propre :
    le cloisonnement des données qu'il agrège, sa sauvegarde et sa confidentialité.
  \item \textbf{Six axes}, lus indépendamment, chacun sur l'intégralité du périmètre qui le
    concerne : contrôle d'accès et autorisations ; échappement des sorties, SQL et portabilité ;
    circuit du stage, calculs et intégrité des données ; imports, exports et PDF ; RGPD,
    sauvegarde, services web et courriels ; tests, internationalisation, documentation et
    maintenabilité.
  \item \textbf{Contrôles automatiques.} \code{php -l} sur les 177 fichiers, \code{phpcs} avec
    le standard \code{moodle} (\code{moodlehq/moodle-cs}) sur la configuration du dépôt, et un
    script comparant les fichiers de langue entre eux et avec les appels du code.
  \item \textbf{Règle de preuve.} Chaque constat est vérifié par lecture du code appelant et des
    fonctions appelées ; un défaut qu'un helper corrige déjà n'est pas retenu. Les lignes citées
    sont celles du commit audité. Rien n'a été corrigé : ce rapport propose, la correction est un
    second temps.
\end{itemize}

\noindent\textbf{Échelle.} \crit~critique : exposition ou altération de données par un acteur
non autorisé, ou perte de données, en usage normal ; \maj~majeure : donnée fausse ou perdue,
décision faussée ou page en erreur dans un cas d'usage réel, mais circonscrit ; \mineure~mineure :
défense en profondeur, cohérence, maintenabilité.

\section{Synthèse}

\textbf{Aucun constat critique.} Les contrôles d'accès sont présents sur toutes les pages, aucune
injection SQL ni d'en-tête de courriel n'a été trouvée, la syntaxe et le standard de codage sont
propres (0 erreur \code{phpcs}), les fichiers de langue sont parfaitement synchronisés (865 clés,
0 manquante, 0 utilisée sans définition), les 19 tables sont sauvegardées et le fournisseur de
confidentialité couvre toutes les colonnes d'utilisateur. La suite de tests compte 226 tests
PHPUnit et 9 scénarios Behat, joués sur trois combinaisons Moodle/PHP/SGBD.

Les constats majeurs se concentrent sur trois familles : des valeurs importées de fichiers qui
ne sont ni nettoyées ni échappées (une XSS stockée), des cas du circuit du stage où un état ou
un calcul reste incohérent (réinitialisation, validation en masse, transfert, fuseau horaire,
durées par année) et des imports qui perdent ou écrasent des données sans le dire.

\begin{center}
\small
\begin{tabularx}{\linewidth}{@{}l X l@{}}
\toprule
& \textbf{Constat} & \textbf{Axe} \\
\midrule
\maj & \og Structure\fg{} importée affichée sans échappement : XSS stockée exécutable chez
  l'étudiant, le référent et la DEVE & Échappement \\
\maj & Réinitialisation d'un stage : l'évaluation ou le motif de refus de l'enseignant survit &
  Circuit \\
\maj & Validation DEVE en masse : les stages refusés sont cochables et validés avec la durée
  déclarée & Circuit \\
\maj & Jours effectifs recalculés dans le fuseau du lecteur : effacés par une édition faite dans
  un autre fuseau & Circuit \\
\maj & Thématique bornée à plusieurs années : seule la durée de la dernière année compte ;
  bilan et tableau de bord divergent & Circuit \\
\maj & Transfert entre promotions : la mobilité d'études n'est pas transférée & Circuit \\
\maj & Champs texte sans limite de longueur vers des colonnes \code{char(255)} : erreur SQL &
  Circuit \\
\maj & Import historique : enseignant évaluateur jamais enregistré & Imports \\
\maj & Import CSV : première ligne perdue sans en-tête ; lignes incomplètes écartées en silence &
  Imports \\
\maj & Réimport StageVet : écrase une évaluation saisie dans Moodle sur un stage SignVet &
  Imports \\
\maj & Gabarit PDF illisible par FPDI : échec sans message, fichiers temporaires abandonnés &
  Imports \\
\maj & Six pages d'import lisent \code{$_FILES} directement : antivirus et quotas de Moodle
  contournés & Qualité \\
\maj & Duplication de cours sans utilisateurs : responsables de thématique perdus & Sauvegarde \\
\maj & Registre RGPD incomplet (5 colonnes d'utilisateur, représentant de l'organisme) ; jours,
  référents et responsables déclarés mais jamais exportés & RGPD \\
\maj & Suivi multi-promotions : une inscription suspendue garde la vue sur ses anciens
  étudiants & Accès \\
\maj & Adresse du maître de stage saisie par l'étudiant, invitation automatique : aucun
  garde-fou technique & Accès \\
\maj & Aucun test sur \code{export.php}, \code{get_my_stages}, l'accès par jeton aux objectifs,
  six tables de la restauration & Tests \\
\maj & Journal des versions arrêté à la 0.5.1 du 3~octobre ; trois points de sauvegarde livrés
  sous le même numéro & Documentation \\
\bottomrule
\end{tabularx}
\end{center}

Les mineures (une cinquantaine) sont détaillées axe par axe. Les plus utiles à traiter avec les
majeures : comptes rendus d'import non échappés, entités HTML dans les PDF, classeurs et
courriels, langue du destinataire ignorée dans les courriels, jeton du maître de stage sans
expiration, 40 chaînes de langue orphelines.

\section{Contrôle d'accès et autorisations}

\subsection{Ce qui tient}

Chaque page fait \code{require_login($course, true, $cm)} puis exige la capacité propre à
l'action : \code{submit} pour l'étudiant, \code{evaluateteacher} pour le référent,
\code{validatedeve} ou \code{registerstages} pour la DEVE, \code{managethemes},
\code{manageteachers}, \code{viewall} pour les exports. Un utilisateur qui n'a que
\code{mod/stage:view} n'atteint aucune donnée d'un autre. Chaque identifiant de saisie, de
thématique, de gabarit ou de liste est relu avec le filtre \code{stageid = $stage->id} : pas
d'accès croisé entre activités. Les fichiers sont cherchés par hachage \emph{dans} la zone de la
saisie ou de la thématique attendue, jamais par hachage global. L'étudiant n'agit que sur ses
saisies ; le référent passe partout par \code{stage_get_assigned_students()}, et une liste vide
donne une page vide, pas tout. Toutes les écritures exigent le \code{sesskey}. \code{returnurl}
est partout en \code{PARAM_LOCALURL}. Les services web font \code{validate_context} puis
\code{require_capability}, et le service est désactivé par défaut. Le jeton du maître de stage
fait 256~bits (\code{random_bytes(32)}), sous index unique, créé par écriture conditionnelle,
révoqué à l'annulation, absent des événements et de la sauvegarde ; il ne donne accès qu'aux
objectifs de la thématique de sa saisie (\code{theme_objective_file.php}, lignes 45--49).

\subsection{Constats}

\begin{itemize}
  \item[\maj] \textbf{A1. Suivi multi-promotions : capacité sans inscription active.}
    \code{mod/stagesynthesis/locallib.php}, lignes 158--200 : le périmètre d'un enseignant est
    recalculé à chaque affichage depuis l'activité source, mais par \code{has_capability()} seul.
    Une inscription suspendue ou expirée conserve ses attributions de rôle : la capacité reste
    vraie alors que \code{require_login()} refuserait l'accès à la source. Scénario : référent dont
    l'inscription à la promotion N$-$1 a été suspendue en fin d'année ; depuis la synthèse d'un
    autre cours, il voit encore le pilotage, la liste à évaluer et le détail de ses anciens
    étudiants, alors que \code{teacher.php} de la source lui est fermé. Correction : exiger aussi
    \code{is_enrolled($context, $userid, '', true)} et la visibilité du module
    (\code{info_module::is_user_visible}).
  \item[\maj] \textbf{A2. Identité du maître de stage.} L'adresse \code{tutoremail} est saisie par
    l'étudiant (\code{student_register_form.php}, lignes 221--226 ;
    \code{convention_request_form.php}, 176--181) et l'invitation à jeton part automatiquement
    (\code{locallib.php}, 6753--6780). Le référent et la DEVE peuvent la corriger à la revue, mais
    rien ne l'impose. Un étudiant qui donne sa propre adresse rédige l'évaluation du maître de
    stage, définitive sauf réinitialisation. Correction : à la revue, signaler une adresse
    appartenant à un compte du cours ou identique à celle de l'étudiant.
  \item[\mineure] \textbf{A3. Relance du maître de stage ouverte à \code{viewall}.}
    \code{entrydetail.php}, lignes 49 et 87--95 : \code{resendtutor} crée le jeton et envoie un
    courriel sous une capacité déclarée en lecture, là où \code{deve.php} exige
    \code{validatedeve}.
  \item[\mineure] \textbf{A4. Jeton sans expiration ni révocation après évaluation.}
    \code{tutor_eval.php} continue d'afficher le nom de l'étudiant, les dates et l'évaluation
    déposée à quiconque détient le lien, indéfiniment ; le jeton transite en GET (journaux,
    Referer). Correction : effacer \code{tutortoken} après \code{stage_apply_tutor_eval()} ou
    n'afficher qu'une confirmation, et refuser un jeton plus vieux qu'un délai.
  \item[\mineure] \textbf{A5. Minimisation.} \code{entrydetail.php}, lignes 54--57 et 177--199 : le
    responsable de thématique, dont le besoin porte sur le stage et les rapports, reçoit la
    convention entière (date de naissance, adresse, téléphone de l'étudiant).
  \item[\mineure] \textbf{A6. Actions à effet de bord par lien GET.} Protégées par \code{sesskey}
    et confirmation JavaScript, donc sans risque CSRF, mais rejouables sans la confirmation
    cliente : \code{register.php:227} (\code{reset}), \code{deve.php:185,193,241},
    \code{entrydetail.php:297}, \code{themes.php:190,221}, \code{evallists.php:85},
    \code{convention_templates.php:141}, \code{transfer.php:161} (transfert irréversible).
    Correction : \code{single_button} POST avec confirmation serveur pour le transfert et les
    suppressions.
  \item[\mineure] \textbf{A7. \code{teacher_assign.php}, lignes 36--37} : un \code{studentid}
    arbitraire est résolu sans \code{is_enrolled()} et son nom complet renvoyé dans l'URL de
    redirection, ce qui permet d'énumérer les noms de tout compte du site.
  \item[\mineure] \textbf{A8. Lectures avant \code{require_login()}.}
    \code{convention_template.php:39-41}, \code{theme_edit.php:41-42},
    \code{evallist_edit.php:44-47} : un visiteur anonyme distingue, par le message d'erreur, un
    identifiant existant d'un inexistant.
  \item[\mineure] \textbf{A9. Service web et thématiques masquées.}
    \code{register_entries.php:86} appelle \code{stage_get_themes()} sans le filtre de visibilité
    que les formulaires DEVE appliquent.
\end{itemize}

\section{Échappement des sorties, SQL et portabilité}

\subsection{Ce qui tient}

Aucune interpolation de valeur utilisateur dans une requête : les tris passent par des cartes
fermées et \code{stage_sql_order()} ne reçoit que des colonnes constantes ; les recherches
utilisent \code{sql_like()} avec \code{sql_like_escape()} ; les listes \code{IN} passent par
\code{get_in_or_equal()} avec des préfixes distincts, découpées par \code{STAGE_SQL_CHUNK} ou
remplacées par une sous-requête ; la pagination est passée en paramètres, jamais dans le texte
SQL ; les \code{GROUP BY} sont complets ; \code{sql_compare_text()} sert à comparer les champs
texte. Rien de propre à MySQL, ce que la CI sur MariaDB confirme. Tous les paramètres sont typés ;
\code{import_preview} range ses fichiers en JSON (pas d'\code{unserialize}), par utilisateur et
par activité, en écriture atomique. L'export Excel n'écrit que des cellules typées
(\code{write_string}) : pas d'injection de formule. Les courriels partent en texte brut par
\code{email_to_user()}, sujets et noms assainis par PHPMailer. Les PDF n'utilisent que
\code{Cell} et \code{MultiCell}, jamais \code{writeHTML}. La page du maître de stage échappe
tout ce qu'elle réaffiche.

\subsection{Constats}

\begin{itemize}
  \item[\maj] \textbf{B1. XSS stockée par la colonne \og structure\fg{} des imports.} Les
    formulaires typent \code{structure} en \code{PARAM_TEXT}, mais aucun importateur ne nettoie
    (\code{csv_importer.php:80} et \code{:737}, \code{historical_importer.php:179},
    \code{import_global.php:216}), et trois affichages rendent la valeur brute :
    \code{locallib.php:4905} (tableau \og Mes stages\fg{} du tableau de bord, \code{html_table}),
    \code{transfer.php:132} et \code{import_stagevet.php:418-421}, où
    \code{stage_entry_short_description()} alimente \code{html_writer::select()}, qui n'échappe
    pas ses libellés. Scénario : un \code{<img src=x onerror=...>} dans un export StageVet, un
    export global ou un CSV remis à la DEVE s'exécute ensuite chez l'étudiant (\code{view.php}),
    les référents et la DEVE (\code{dashboard.php}, synthèse). Correction : \code{s()} aux trois
    affichages et \code{clean_param(..., PARAM_TEXT)} sur chaque champ texte lu d'un fichier dans
    les quatre importateurs.
  \item[\mineure] \textbf{B2. Comptes rendus d'import non échappés.} \code{import.php:86-89},
    \code{teachers_import.php:85-88}, \code{import_stagevet.php:353-361, 455-463, 467-470} :
    courriels, noms et thématiques lus dans le fichier sont injectés dans une
    \code{notification()}, qui ne purifie pas. \code{study_mobility_import.php:106},
    \code{import_global.php:300-303} et \code{import_historical.php:337-340} font correctement
    \code{array_map('s', ...)}.
  \item[\mineure] \textbf{B3. \code{format_string()} vers des sorties non HTML.} Les entités
    \code{&amp;} et \code{&#039;} apparaissent littéralement dans la page~1 de la convention
    (\code{locallib.php:8107}), les classeurs (\code{export.php:347,357,395,509}), les courriels
    texte (\code{locallib.php:2663-2665, 2698-2701, 6581}) et les noms de fichiers
    (\code{export.php:49}, \code{export_pdf.php:128}). Correction :
    \code{format_string($x, true, ['escape' => false])}.
  \item[\mineure] \textbf{B4. Noms de fichiers sans \code{s()}} dans \code{locallib.php:7096} et
    \code{:7798} ; non exploitable (\code{PARAM_FILE} à l'envoi), défense en profondeur.
  \item[\mineure] \textbf{B5. Fichiers temporaires de la convention.} \code{locallib.php:7903}
    écrit dans \code{sys_get_temp_dir()} et non \code{make_temp_directory('mod_stage')} ; les
    \code{unlink} (8194--8196) ne sont pas dans un \code{finally} (voir D5).
\end{itemize}

\section{Circuit du stage, calculs et intégrité des données}

\subsection{Ce qui tient}

Six statuts de stage, huit de convention, et chaque transition passe par les écritures
conditionnelles introduites après le test de concurrence. \code{upgrade.php} est idempotent et
complet par rapport à \code{install.xml}, points de sauvegarde croissants, \code{version.php}
égal au dernier. Un index unique existe partout où le code suppose l'unicité (jeton, triplet
référent, jours, réponses, check-list, durées). \code{stage_delete_instance()} couvre les 19
tables ; \code{stage_delete_entries()} purge plages, jours, détail, réponses, check-list et les
deux zones de fichiers, par paquets ; la suppression d'une thématique, d'un gabarit, d'une liste
ou d'une question est refusée tant que des saisies ou des réponses en dépendent. Les annulés et
non validés sont exclus de tous les décomptes, les bornes de plages sont incluses, l'itération
des jours passe par \code{DateTimeImmutable::modify}, donc sans dérive à l'heure d'été ; les
chevauchements et les plages de plus de 366 jours sont refusés. Les événements ont une classe de
base commune, \code{relateduserid} imposé, et relisent le statut en base avant de se déclencher.
Les tâches planifiées portent leurs marqueurs d'idempotence.

\subsection{Constats}

\begin{itemize}
  \item[\maj] \textbf{C1. Réinitialisation incomplète.} \code{stage_reset_entry()}
    (\code{locallib.php:2355-2380}) remet statut, contournement, décision DEVE et durée retenue à
    zéro, mais laisse \code{teacherid}, \code{teachereval}, \code{teachertime}. Un stage refusé
    par le référent (motif dans \code{teachereval}) puis réinitialisé affiche \og Évalué par X le
    \dots\fg{} avec le motif de refus sur un stage \og Enregistré\fg{} (\code{entrydetail.php:
    243-256}) ; la DEVE revoit cette évaluation périmée dès que le stage repasse en auto-évalué
    (\code{deve.php:146-152}) ; l'export la sort. Si le référent réévalue ensuite par
    questionnaire, l'ancien motif reste pour toujours. Correction : remettre les trois champs à
    \code{null}.
  \item[\maj] \textbf{C2. Validation en masse des stages refusés.} \code{deve.php:268-283} et
    \code{336-345} : la liste est filtrée par \code{status < 3}, donc les \og Non validé\fg{}
    y figurent avec leur case à cocher (seuls les annulés en sont privés). \og Tout
    sélectionner\fg{} puis \og Valider\fg{} les valide avec la durée déclarée et écrase
    \code{devecomment} par une chaîne vide, alors que la validation unitaire propose les jours
    effectifs cochés et que la doc demande une réinitialisation préalable. Correction : exclure
    \code{status < 0} de la sélection et de la boucle ; retenir le nombre de jours cochés s'il
    existe, comme en unitaire.
  \item[\maj] \textbf{C3. Jours effectifs et fuseau horaire.} \code{stage_get_period_days()}
    (\code{locallib.php:2003-2013}) calcule les minuits dans le fuseau de l'utilisateur
    \emph{courant} ; \code{stage_set_entry_workdays()} (2040--2078) ne conserve que les jours
    présents dans ce calcul ; \code{stage_save_entry_periods()} (1795) rejoue
    \code{stage_set_entry_workdays()} avec les jours existants. Un étudiant en stage à l'étranger
    avec son fuseau réglé (par exemple \code{America/Montreal}) coche ses jours ; la DEVE, à
    Paris, modifie la saisie sans toucher aux dates (\code{register.php:306-317}) : aucun jour
    ne correspond plus, tous sont supprimés en silence. La grille des jours se décale aussi d'un
    jour selon le lecteur. Correction : normaliser sur le fuseau du site (ou stocker des
    \code{AAAAMMJJ}) et ne jamais filtrer les jours existants dans un autre fuseau que celui du
    stockage.
  \item[\maj] \textbf{C4. Durées par année d'une thématique bornée.}
    \code{stage_get_student_year_progress()} (\code{locallib.php:745-760}) ne lit, pour une
    thématique A2--A4, que la durée de l'année finale et la compare au cumul ; les durées propres
    à A2 et A3, que \code{theme_edit.php} met pourtant en avant, ne sont jamais lues.
    \code{stage_get_student_progress()} (1013--1035), lui, somme les années où l'étudiant a une
    saisie. Avec 10/10/10 jours et un seul stage de 10 jours en A2, le bilan de fin de cursus
    annonce \og objectifs atteints\fg{} (requis 10) ; le tableau de bord étudiant, dès qu'un
    stage A4 s'ajoute, demande 20. Décision de fin de cursus potentiellement faussée.
    Correction : requis = somme des durées de la plage, ou refuser la saisie par année quand une
    plage est définie.
  \item[\maj] \textbf{C5. Transfert sans la mobilité d'études.}
    \code{stage_execute_student_transfer()} (\code{locallib.php:5375-5445}) déplace saisies,
    réponses, check-lists, fichiers, et supprime les attributions source, mais ignore
    \code{stage_study_mobility} (clé \code{stageid, userid}). Une mobilité validée dans la source
    met l'étudiant \og en défaut\fg{} de jours à l'étranger dans la cible, et la ligne reste
    orpheline. Correction : déplacer la ligne, en tenant compte de l'index unique côté cible.
  \item[\maj] \textbf{C6. Longueur des champs texte.} Aucune règle \code{maxlength} ni troncature
    sur \code{structure} et \code{country} (\code{deve_entry_form.php:78-94},
    \code{student_register_form.php:79-86}, \code{register.php:396-399},
    \code{register_entries.php:45,61}), le nom de thématique (\code{theme_edit_form.php:84}), le
    nom de gabarit (\code{convention_template_form.php:46}) et les champs d'établissement
    (\code{conventions_admin_form.php:75}, \code{establishmentphone} en \code{char(64)}). Une
    raison sociale de 300 caractères collée par un étudiant donne une
    \code{dml_write_exception} : page d'erreur, rien d'enregistré. Seuls
    \code{stage_convention_detail} (longueurs lues du schéma) et les imports CSV sont protégés.
  \item[\mineure] \textbf{C7. Contournement de l'évaluation du maître de stage}
    (\code{locallib.php:6873-6885}) : \code{update_record} inconditionnel, possible sur un stage
    annulé ou validé, sans événement.
  \item[\mineure] \textbf{C8. Redemande de convention après refus} (\code{locallib.php:6192-6208}) :
    \code{conventionteachervalidatedby} et sa date ne sont pas effacés ; la fiche affiche
    \og Validée par X\fg{} sur une demande qui attend de nouveau le référent.
  \item[\mineure] \textbf{C9. Import StageVet sans événement} (\code{csv_importer.php:735-823}) :
    insertion directe en statut auto-évalué avec \code{tutortime}, sans \code{entry_created} ni
    les événements d'évaluation ; les journaux ne retrouvent pas ces stages.
  \item[\mineure] \textbf{C10. Édition DEVE d'une saisie} (\code{locallib.php:1638-1684},
    \code{register.php:280-317}) : \code{update_record} de l'objet entier lu en début de page,
    aucun événement pour un changement de thématique, d'année ou de durée après validation, et
    \code{declaredduration} préremplie par le nombre de jours cochés : toute édition, même d'un
    autre champ, remplace la durée déclarée.
  \item[\mineure] \textbf{C11. Invitations du maître de stage}
    (\code{send_tutor_evaluation_requests.php:95-125}, \code{locallib.php:6649-6695}) : ni
    \code{cm.visible} ni le compte étudiant (suspendu, désinscrit, supprimé) ne sont testés, alors
    que la relance de convention le fait (\code{send_convention_reminders.php:85-94}).
  \item[\mineure] \textbf{C12. Durée retenue} (\code{locallib.php:2231-2256}) : ni bornée par les
    plages, ni possible à zéro (\code{retainedduration > 0 ? ... : declaredduration}).
  \item[\mineure] \textbf{C13. Relance de convention non rejouée} après report des plages :
    \code{conventionremindertime} n'est jamais remis à \code{null} (\code{locallib.php:6476-6515}).
  \item[\mineure] \textbf{C14. Levée de dispense} possible sur un stage déjà évalué
    (\code{locallib.php:6160-6190}).
  \item[\mineure] \textbf{C15. Divers schéma et \code{lib.php}} : \code{install.xml} porte
    \code{VERSION="20261026"} pour un dernier point \code{2026103101} ;
    \code{stage_question.themeid} est \code{NOT NULL} avec clé et index mais toujours créé à 0
    (\code{evallist_edit.php:198}) ; \code{theme_edit.php:184-186} accepte une durée négative ;
    \code{stage_user_outline()} renvoie \code{time()} ; \code{MOD_PURPOSE_COMMUNICATION} pour une
    activité d'évaluation.
\end{itemize}

\section{Imports, exports et PDF}

\subsection{Ce qui tient}

Tous les imports écrivent sous le verrou de l'activité et en transaction, avec relecture des
doublons sous le verrou : jamais d'import à moitié appliqué. Les en-têtes sont normalisés
(casse, accents, apostrophes, BOM), les libellés résolus dans toutes les langues installées, les
homonymes jamais tranchés par le code (correspondance unique ou arbitrage DEVE), les thématiques
inconnues toujours refusées, les dates strictes (\code{checkdate}, refus des numéros de série
Excel et des plages inversées). StageVet n'écarte aucune ligne en silence, détecte doublons et
chevauchements intra-fichier, persiste les arbitrages en JSON et rejoue le même contenu. Les
prévisualisations sont en JSON, par utilisateur, atomiques, périmées après un jour. Le XLSX est
lu avec \code{LIBXML_NONET}. L'export Excel lit chaque table rattachée en une requête. Les
importateurs ont des tests réels, et le PDF de convention est testé avec un vrai gabarit FPDI.

\subsection{Constats}

\begin{itemize}
  \item[\maj] \textbf{D1. Import historique : enseignant évaluateur perdu.}
    \code{import_historical.php:235-237} pose \code{teacherid}, \code{teachertime},
    \code{teachereval} sur l'objet, puis \code{stage_apply_deve_validation()} n'écrit que
    \code{deveuserid}, \code{devetime}, \code{retainedduration}, \code{status},
    \code{timemodified} et \code{devecomment} (\code{locallib.php:2231-2260}). L'enseignant est
    bien ajouté comme référent, mais aucun stage historique n'a d'évaluateur : \og Évalué par\fg{}
    vide dans l'export global. Correction d'une ligne : \code{stage_update_entry_fields($entry,
    ['teacherid', 'teachertime', 'teachereval'])} dans la même transaction.
  \item[\maj] \textbf{D2. Import CSV simple : en-tête et lignes incomplètes.}
    \code{csv_importer.php:58-65} (et \code{:221} pour les référents) consomme toujours la
    première ligne comme en-tête, alors que l'aide (\code{lang/fr/stage.php:115,220}) l'annonce
    facultative : sans en-tête, le premier stage n'est ni créé ni signalé. Ligne 87, une ligne
    dont le courriel ou la thématique est vide est sautée sans erreur, à l'inverse de la règle de
    StageVet (522--534). Dans les deux cas le compte rendu annonce \og N importés\fg{} sans dire
    ce qui manque.
  \item[\maj] \textbf{D3. Réimport StageVet sur un stage SignVet.}
    \code{csv_importer.php:1046-1064} : pour un stage \code{SIGNVET}, toute évaluation non vide de
    l'export remplace l'existante quel que soit le statut, et \code{tutortime} est remis à
    \code{time()} ; la protection \og conservée\fg{} ne joue que pour les autres stages, et le test
    (\code{stagevet_evaluations_test.php:236}) ne couvre que ceux-là. Une auto-évaluation saisie
    dans Moodle ou une réponse du maître de stage par jeton est écrasée au réimport suivant, sans
    ligne de rapport.
  \item[\maj] \textbf{D4. Gabarit PDF illisible.} \code{locallib.php:8186} :
    \code{setSourceFile()} lève une exception FPDI pour un PDF 1.5+ à table de références
    compressée (Word, Acrobat \og optimisé\fg{}) ou chiffré, sans \code{try/catch} ;
    \code{convention_template.php} ne tente pas de lire le gabarit à l'envoi. Dans
    \code{convention.php:143-149} le téléchargement se fait dans un cadre invisible : la DEVE voit
    \og Téléchargement en cours\fg{} puis rien ; \code{convention_review.php} passe la demande en
    \og éditée\fg{} sans que le PDF ait été construit. Correction : ouvrir le gabarit avec FPDI à
    l'enregistrement et refuser le formulaire avec un message ; en génération, encapsuler et
    renvoyer une erreur lisible.
  \item[\mineure] \textbf{D5. Fuite de fichiers temporaires} : les copies du gabarit et des logos
    ne sont supprimées qu'en cas de succès (\code{locallib.php:8194}) ; chaque échec D4 en laisse
    trois sur disque, hors du répertoire temporaire de Moodle (B5).
  \item[\mineure] \textbf{D6. Analyse CSV} : \code{declaredduration = (int) trim(...)}
    (\code{csv_importer.php:83}) accepte \og 12,5\fg{}, \og abc\fg{} ou \og -3\fg{} sans
    signalement ; UTF-16 non détecté (\code{detect_encoding}, 956) ; \code{parse_date} (1121)
    produit minuit dans le fuseau de l'importateur, si bien que deux membres de la DEVE de fuseaux
    différents ne se voient pas en doublon ; fichier chargé entièrement en mémoire sans
    \code{raise_memory_limit} (\code{import.php:66}, \code{import_stagevet.php:207},
    \code{teachers_import.php:65}) ; lignes StageVet à thématique inconnue non reprises dans les
    lignes en attente (585) ; \og 0 stage importé\fg{} affiché après un arbitrage vide (229).
  \item[\mineure] \textbf{D7. Prévisualisations} : \code{import_preview.php:71} écrit avec le
    umask par défaut et non \code{$CFG->filepermissions} ; le formulaire de confirmation ne porte
    pas d'identifiant de prévisualisation (deux onglets du même utilisateur confirment le dernier
    classeur analysé) ; \code{import_historical.php:172} supprime la prévisualisation avant la
    transaction.
  \item[\mineure] \textbf{D8. Lecteur XLSX} (\code{historical_importer.php}) : \code{getFromName()}
    renvoyant \code{false} donne une \code{ValueError} (403) non attrapée dans
    \code{study_mobility_import.php:90}, qui n'attrape que \code{moodle_exception} : page d'erreur
    fatale pour un \code{.docx} renommé ; attribut \code{r} absent, lignes écrasées sur l'index
    $-1$ (446) ; colonne \og année\fg{} par défaut (221) ; empreinte de doublon sans les dates
    (\code{import_historical.php:118-121}).
  \item[\mineure] \textbf{D9. Restauration d'un export global} : deux requêtes par ligne à la
    prévisualisation (\code{import_global.php:178,266}, soit 60\,000 pour 30\,000 stages) ;
    exportés mais jamais restaurés : référents de l'étudiant (\code{export.php:436}), responsables
    de thématique, mode de rapport, feuilles 1 et 3 ; jamais exportés : jours effectifs,
    check-lists, réponses, fichiers, mobilité d'études ; année \og non précisée\fg{} absente des
    options (144, un avertissement par ligne) ; oui/non reconnus en français et anglais seulement
    (55) ; aucune détection de doublon à l'intérieur du classeur (52) ; auto-évaluation convertie
    en texte (\code{export.php:435}) puis réaffichée en \code{FORMAT_HTML}.
  \item[\mineure] \textbf{D10. PDF} : \code{'VetAgro Sup'} en dur (\code{convention_pdf.php:96}) ;
    logo \code{.png} non vérifié (un JPEG renommé fait lever TCPDF au moment de la génération, même
    chemin invisible que D4) ; autoloader Composer de \code{thirdparty/vendor} non déclaré dans
    \code{thirdpartylibs.xml}.
  \item[\mineure] \textbf{D11. Zip des rapports} (\code{locallib.php:7876}) : deux fichiers d'une
    même saisie dont les noms nettoyés coïncident, le second remplace le premier.
  \item[\mineure] \textbf{D12. Requêtes par ligne} : \code{stage_get_signed_convention_file()} par
    saisie dans \code{register.php:215}, \code{conventions.php:136}, \code{teacher.php:309} ;
    \code{stage_get_answers()} relu quatre fois dans \code{teacher.php} ; une requête par étudiant
    pour le modèle de mobilité (\code{study_mobility_import.php:67}).
\end{itemize}

\section{RGPD, sauvegarde, services web et courriels}

\subsection{Ce qui tient}

\code{get_contexts_for_userid()} et \code{get_users_in_context()} couvrent toutes les colonnes
d'utilisateur du schéma, vérifiées une à une. La suppression distingue l'étudiant (tout est
effacé, fichiers compris) du personnel (dissociation à \code{null} et effacement de ses textes),
avec des colonnes déclarées nullables. La sauvegarde embarque les 19 tables, respecte
\code{userinfo}, annote toutes les colonnes d'utilisateur, exclut volontairement le jeton (pas de
collision à la restauration sur le même site) et n'applique aucun décalage de dates. Les
services web valident contexte et capacité, et \code{get_my_stages} ne renvoie que le strict
nécessaire. Les courriels partent de l'adresse sans réponse, avec des liens absolus, un
utilisateur factice complet pour le maître de stage et \code{validate_email()} avant envoi.
\code{mod_stagesynthesis} est un \code{null_provider} exact, et remappe ses liens à la
restauration, avec test.

\subsection{Constats}

\begin{itemize}
  \item[\maj] \textbf{E1. Registre des traitements incomplet.} \code{provider.php:69-89} ne
    déclare pas \code{conventionrejectedby}, \code{conventionteachervalidatedby},
    \code{conventioneditedby}, \code{conventionsignedby}, \code{cancelledby}
    (\code{install.xml:216-221}), ni \code{conventionrejectcomment}, \code{tutortoken} et les
    dates d'étapes ; lignes 91--101, les données nominatives du représentant de l'organisme
    d'accueil (\code{hostrepresentative}, \code{hostphone}, \code{hostemail},
    \code{hostaddress}...) ne sont pas déclarées. Le registre généré par Moodle omet qu'un membre
    de la DEVE est nommément enregistré comme signataire ou annulateur.
  \item[\maj] \textbf{E2. Tables déclarées mais jamais exportées.} \code{stage_entry_workday}
    (métadonnées 108--110, aucune lecture dans \code{export_user_data}, 333--417),
    \code{stage_entry_teacher} et \code{stage_theme_teacher} (ni côté étudiant, ni dans
    \code{export_staff_data}, 443--505), \code{stage_study_mobility.validatedby} côté DEVE. Un
    étudiant dont le stage est compté jour par jour reçoit un export sans ses jours.
  \item[\maj] \textbf{E3. Dissociation du personnel incomplète.} \code{teacher.php:70,104}
    enregistre l'évaluation du référent dans \code{stage_answer} (questions \code{teacher}),
    table sans colonne d'auteur ; à la suppression de l'enseignant, \code{provider.php:627-652}
    efface \code{teachereval} mais laisse ces réponses, contrairement à la politique énoncée aux
    lignes 53--56. À trancher : les effacer, ou documenter qu'elles appartiennent au dossier de
    l'étudiant.
  \item[\maj] \textbf{E4. Duplication sans utilisateurs : responsables de thématique perdus.}
    \code{stage_theme_teacher} est sauvegardé sans condition
    (\code{backup_stage_stepslib.php:233-236}, commenté comme du paramétrage), mais
    \code{process_stage_theme_teacher()} (\code{restore_stage_stepslib.php:177-180}) exige
    \code{get_mappingid('user')}, vide quand les utilisateurs ne sont pas dans l'archive
    (duplication d'activité, import de cours). La DEVE qui duplique l'activité pour la promotion
    suivante ressaisit tous les responsables. Correction : sur le même site, conserver l'id
    d'origine si le compte existe, comme le fait la restauration de \code{mod_stagesynthesis}.
  \item[\mineure] \textbf{E5. Export étudiant partiel} (\code{provider.php:336-365}) :
    \code{conventionrejectcomment}, \code{abroad}, dates d'évaluation et de convention,
    \code{tutorbypassed} absents ; \code{referentteacherid} exporté comme entier et non comme nom.
  \item[\mineure] \textbf{E6. Emplacement externe non déclaré} pour le courriel au maître de stage
    (nom complet de l'étudiant et lien à jeton envoyés hors Moodle) :
    \code{add_external_location_link}.
  \item[\mineure] \textbf{E7. Noms de personnes dans le paramétrage} :
    \code{establishmentrepresentative}, \code{establishmentsignatory} (\code{install.xml:17,21}).
  \item[\mineure] \textbf{E8. Tests de confidentialité et de restauration} : pas d'export de
    l'étudiant, pas de jour, check-list ni fichier créés avant suppression, pas de
    \code{provider_testcase} ; \code{backup_restore_test.php} n'affirme que 12 tables sur 19
    (manquent \code{stage_theme_duration}, \code{stage_entry_workday},
    \code{stage_convention_detail}, \code{stage_study_mobility},
    \code{stage_evallist_question}, \code{stage_email_template}), ni restauration sans
    utilisateurs (E4), ni utilisateur absent.
  \item[\mineure] \textbf{E9. Décodage des liens} : \code{stage_theme_checklist.description} et
    \code{leavemodalities} absents de \code{define_decode_contents}.
  \item[\mineure] \textbf{E10. Services web} : classes globales \code{\external_api} et
    \code{externallib.php}, dépréciées depuis Moodle 4.2 (MDL-76583) pour un module qui exige
    4.3 ; \code{register_entries.php:273-275} écarte en silence un étudiant non inscrit ou une
    thématique invalide, sans les nommer dans la réponse ; \code{readonlysession} absent sur
    \code{get_my_stages} ; \code{get_my_stages} sans test.
  \item[\mineure] \textbf{E11. Courriels} : \code{stage_resolve_email_text()} appelé sans
    \code{$lang} (\code{locallib.php:2661, 2695, 6250, 6423, 6456, 6537}), donc dans la langue
    de l'acteur ou du cron, pas du destinataire ; pas de \code{db/messages.php} (préférences de
    notification et mode numérique contournés) ; corps personnalisé en \code{PARAM_RAW}
    (\code{notifications_form.php:69-70}) envoyé en texte brut ; étudiant lu sans
    \code{deleted = 0} (\code{locallib.php:6416, 6449}), d'où un \code{debugging()}.
\end{itemize}

\section{Tests, internationalisation, documentation et maintenabilité}

\subsection{Ce qui tient}

Fichiers de langue parfaitement synchronisés : 865 clés de chaque côté, 0 manquante, 0 utilisée
sans définition, 0 décalage de \code{{$a}} ; aucun texte littéral dans les gabarits Mustache.
Aucune API Moodle dépréciée, aucune syntaxe PHP postérieure à 8.0 (la promesse de
\code{version.php} est tenue et vérifiée par la matrice CI), PHPDoc sur 100~\% des fonctions,
contrôlée en CI de façon bloquante. CI exemplaire pour un plugin : \code{phpcs} avec
avertissements bloquants, bornes basse et haute de Moodle, PostgreSQL et MariaDB, \code{phplint},
\code{validate}, \code{savepoints}, \code{phpdoc}, \code{mustache}, PHPUnit pour les deux
modules, Behat sans Selenium, \code{persist-credentials: false}, permissions minimales, et un
commentaire justifiant chaque choix. 226 tests dans 38 classes, 162 \code{@covers}, tests de
régression dédiés (concurrence, audit, circuit), générateur complet. \code{progress_cache} est un
cache de requête, armé et vidé dans un \code{finally}, sans invalidation à gérer.
\code{thirdpartylibs.xml} déclare exactement FPDI 2.6.8.

\subsection{Constats}

\begin{itemize}
  \item[\maj] \textbf{F1. Envois de fichiers hors API.} Six pages lisent \code{$_FILES} puis
    \code{file_get_contents($tmp_name)} (\code{import.php:58}, \code{import_global.php:97},
    \code{import_historical.php:293}, \code{import_stagevet.php:202},
    \code{study_mobility_import.php:83}, \code{teachers_import.php:57}). Le contrôle
    \code{UPLOAD_ERR_OK} + \code{is_uploaded_file()} est correct, mais ce chemin contourne l'API
    fichiers de Moodle, donc l'analyse antivirus (\code{core\antivirus}) et les quotas. Correction :
    élément \code{filepicker} de \code{moodleform} et \code{get_file_content()}, factorisés.
  \item[\maj] \textbf{F2. Zones à risque sans test.} \code{export.php} (566 lignes, export
    réglementaire) : 0 référence ; \code{promotion_pdf} : 0 ; \code{get_my_stages} : 0 ;
    \code{theme_objective_file.php}, seul accès à un fichier sans \code{require_login} : 0 ; six
    tables de la restauration (E8) ; tri, pagination et filtres des listes (\code{stage_sql_order},
    \code{stage_count_filtered_entries}, \code{stage_get_filtered_entries_page}), c'est-à-dire
    le SQL écrit à la main qui justifie le job MariaDB : aucun test direct. Au total 107 des 267
    fonctions de \code{locallib.php} et \code{lib.php} (40~\%) et 30 des 55 pages ne sont nommées
    dans aucun test. Behat couvre le circuit nominal, les questions notées et l'auto-enregistrement,
    mais ni import, export, convention PDF, maître de stage, suppression en masse ni transfert.
  \item[\maj] \textbf{F3. Journal des versions.} La dernière section versionnée de
    \code{CHANGELOG.md} est \og mod\_stage 0.5.1 --- 2026-10-03\fg{} (ligne 225) ; depuis,
    trois points de sauvegarde (\code{2026102600}, \code{2026103100}, \code{2026103101}) et
    quatre commits sur \code{version.php} ont été livrés sous la même chaîne \code{0.5.1},
    pendant que \og Non publié\fg{} accumule 200 lignes (mobilité d'études, concurrence,
    performance, Moodle 4.3 minimum). Le fichier contredit sa propre règle (lignes 4--6). De même,
    \code{mod_stagesynthesis} annonce un nouvel événement et un changement de clé de filtre sans
    changer de version (\code{2026100200}), et les guides portent \code{\versionstage{0.5.1}}.
    Correction : publier 0.6.0, vider \og Non publié\fg{}, incrémenter la synthèse.
  \item[\mineure] \textbf{F4. Chaînes de langue} : 40 clés définies et jamais utilisées, résidus
    des pages refondues en 0.5.0 (\code{addquestion}, \code{questionsaved},
    \code{managethemedurations}, \code{availableteachers}, \code{addchecklistitem}...) ; 4
    chaînes avec HTML dont \code{deleteentrieswarning_export} (\code{lang/en/stage.php:1101}) qui
    mêle balisage et texte ; dix \code{get_string(...) . ' : ' . $valeur} (typographie française
    en anglais : \code{entrydetail.php:248,270,284,312}, \code{locallib.php:380,3566,4795},
    \code{theme_edit.php:262}, \code{evallist_edit.php:222}, \code{convention_template.php:106}) ;
    les 17 \code{get_description()} des événements sont en français alors que Moodle attend de
    l'anglais non traduit ; nom du service web en français (\code{services.php:48}) ; langues des
    gabarits et questions limitées à fr/en par le modèle de données (colonne \code{nameen}) ;
    \og DEVE\fg{} 21 fois dans la langue anglaise.
  \item[\mineure] \textbf{F5. Structure du code} : \code{locallib.php} tient 26 groupes
    fonctionnels identifiables (statuts, thématiques et progression, attributions, verrou et
    transitions, périodes et jours, circuit, courriels, listes et questions, rendu du détail,
    filtres et pagination, pilotage, tableau de bord, transfert et import, conventions, relances,
    maître de stage, fichiers, check-lists, responsables, zip, PDF, suppression) sans un seul
    séparateur ; \code{stage_print_student_dashboard_content} fait 356 lignes,
    \code{stage_plan_student_transfer} 197, \code{stage_build_convention_pdf} 175 ; le prologue
    de page (\code{required_param}, \code{get_coursemodule_from_id}, \code{get_course},
    \code{get_record}, \code{require_login}, contexte) est copié dans 50 pages ; 680 appels
    \code{html_writer} contre 6 \code{render_from_template}.
  \item[\mineure] \textbf{F6. CI et dépôt} : couverture mesurée mais sans seuil ; aucun
    \code{export-ignore} dans \code{.gitattributes} (une archive embarque tests, docs et CI) ;
    \code{.gitignore} sans les auxiliaires LaTeX ; actions épinglées par tag et non par
    empreinte ; pas de détection de code dupliqué.
  \item[\mineure] \textbf{F7. Documentation} : \code{INSTALL.md} décrit quatre rôles sans
    tableau capacité par archétype ; la tâche de relance des conventions (6~h~30,
    \code{db/tasks.php}) n'est documentée nulle part, seules les invitations de 6~h~45 le sont
    (\code{INSTALL.md:401-402}, \code{README.md:70}).
\end{itemize}

\section{Ordre de correction proposé}

\begin{enumerate}[itemsep=2pt]
  \item \textbf{Sécurité et données, un correctif chacun} : B1 (\code{s()} et
    \code{clean_param} dans les importateurs), B2, D1 (une ligne), C1, C2, C5, F1 (filepicker).
  \item \textbf{Calculs qui faussent une décision} : C4 (durées par année d'une plage), C3
    (fuseau des jours effectifs), C6 (\code{maxlength}), D3 (réimport StageVet).
  \item \textbf{Robustesse des conventions} : D4 et D5 (gabarit vérifié à l'envoi,
    \code{try/finally}, message d'erreur), D10 (logo), B3 (entités dans PDF et courriels).
  \item \textbf{RGPD et sauvegarde} : E1, E2, E3 (ou sa documentation), E4, E6 ; tests E8.
  \item \textbf{Accès} : A1 (inscription active dans la synthèse), A2 (signal à la revue), A3,
    A4 (expiration du jeton), A6 (POST pour le transfert et les suppressions).
  \item \textbf{Tests} : F2, dans l'ordre \code{theme_objective_file} par jeton,
    \code{export.php}, \code{get_my_stages}, restauration des six tables, tri et pagination.
  \item \textbf{Publication} : F3 (version 0.6.0, journal, guides, synthèse), F4 (chaînes
    orphelines, descriptions d'événements), F7.
  \item \textbf{Fond de maintenabilité}, au fil des chantiers : F5 (découpage de
    \code{locallib.php}, prologue de page), E10 (\code{core_external}), E11
    (\code{db/messages.php}, langue du destinataire), D9, D12.
\end{enumerate}

\end{document}
